\chapter{MongoDB Indexing}

\section{What an Index Actually Does}

Without an index, MongoDB has to perform a \textbf{collection scan} --- check
every single document against the query filter, one by one. An index is a
separate, ordered data structure (conceptually similar to a book's index)
that maps field values directly to the documents that contain them, so
MongoDB can jump straight to matches instead of scanning everything.

\begin{diagrambox}[Without an Index]
\begin{verbatim}
Query: { status: "completed" }
        |
        v
   Scan document 1 -> check status -> no match
   Scan document 2 -> check status -> no match
   Scan document 3 -> check status -> match!
   ... (continues for every document in the collection)
\end{verbatim}
\end{diagrambox}

\begin{diagrambox}[With an Index]
\begin{verbatim}
Query: { status: "completed" }
        |
        v
   Look up "completed" in the status index
        |
        v
   Index points directly at matching documents
        |
        v
   Return matches (no scan of unrelated documents)
\end{verbatim}
\end{diagrambox}

\section{Unique Indexes}

A unique index both speeds up lookups \emph{and} enforces that no two
documents can share the same value for that field. Declaring
\texttt{unique: true} on a schema field creates this index automatically.

\begin{lstlisting}[language=JavaScript]
// backend/src/models/User.js
const userSchema = new mongoose.Schema({
  email: {
    type: String,
    required: true,
    unique: true, // creates a unique index on `email`
    lowercase: true,
    trim: true,
  },
  password: { type: String, required: true },
});
\end{lstlisting}

\begin{notebox}[NOTE]
A unique index is what actually causes the duplicate-email
\texttt{E11000 duplicate key error} at the database layer, which your error
handler can translate into a clean \texttt{409 Conflict} response (see
Chapter~25).
\end{notebox}

\section{Compound Indexes}

A compound index covers multiple fields together, and is especially valuable
when a query filters on one field and sorts on another --- exactly the
search/filter/sort pattern from Chapter~28.

\begin{lstlisting}[language=JavaScript]
// backend/src/models/Task.js
const taskSchema = new mongoose.Schema({
  title: { type: String, required: true },
  status: { type: String, enum: ["pending", "in-progress", "completed"] },
  owner: { type: mongoose.Schema.Types.ObjectId, ref: "User" },
}, { timestamps: true });

// Speeds up: Task.find({ status: "completed" }).sort({ createdAt: -1 })
taskSchema.index({ status: 1, createdAt: -1 });

module.exports = mongoose.model("Task", taskSchema);
\end{lstlisting}

The order of fields in a compound index matters: this index efficiently
supports queries that filter on \texttt{status} alone, or on
\texttt{status} + sort by \texttt{createdAt}, but is less useful for a query
that sorts by \texttt{createdAt} without filtering on \texttt{status} at all.

\section{When to Add an Index}

\begin{notebox}[NOTE]
Add an index when a field is frequently used to \textbf{filter}
(\texttt{status}, \texttt{owner}), \textbf{sort} (\texttt{createdAt}), or
\textbf{look up uniquely} (\texttt{email}). These are exactly the fields
that appear in \texttt{.find()}, \texttt{.sort()}, or lookup queries across
the app's most common requests.
\end{notebox}

\begin{warnbox}[WARNING]
Indexes are not free. Every index has to be updated on every insert, update,
and delete, which slows down writes, and every index consumes additional
disk and memory. Don't index every field ``just in case'' --- only index
fields that are actually queried or sorted on frequently, and periodically
review indexes on collections that see heavy write traffic.
\end{warnbox}
